\subsection{Progress beyond correct confirmation}\label{sec:interpretation}
For fresh independent private draws, Equation~\eqref{eq:progress} describes
the exact anonymous rule's probability of writing. Proposition~\ref{prop:diffuse}
makes the missing hypotheses explicit, including a bound on each neighbor's
per-color offer probability. This covers the exact always-active Johansson
variant, whereas coincident deterministic proposals can make progress zero
(Counterexample~\ref{ex:stalled}). A finite draw probe demonstrates observed
concentration, not exactly zero probability of an alternative color.

The GM suffix condition in Section~\ref{sec:relevance} explains why a correct
veto can leave an unresolved conflict without implying that every omitted
repair must fail. The two candidate interventions distinguish these
mechanisms in specific copied executions. The broader classification
extends their observed scope without converting a retrospective replay into
a guarantee over new graphs.

The finite-type MAX condition in Appendix~\ref{app:math} certifies one
canonical COMMIT predicate when all 64 type subsets satisfy its sign
inequalities. It does not certify candidate diversity, scheduled repair or
palette-context invariance. All 45 selected canonical transfer cores fail
that sufficient condition; five different historical E1 weights pass
conditional certificates. Successful finite runs do not establish its
hypotheses, and failure of a sufficient certificate is not itself a
reachable failure. The identifier obstruction and representation controls
remain separate from numerical sensitivity (Appendix~\ref{app:ids}).
